% This LaTeX document needs to be compiled with XeLaTeX.
\documentclass[
  10
]{article}
\usepackage[margin=2cm,includehead=true,includefoot=true,centering,]{geometry}
\usepackage{xcolor}
\usepackage{ucharclasses}
\usepackage{hyperref}
\usepackage{amsmath,amssymb}
\usepackage{amsfonts}
\usepackage[version=4]{mhchem}
\usepackage{stmaryrd}
\usepackage{bbold}
\usepackage{polyglossia}
\usepackage{fontspec}
\usepackage{ctex}
\usepackage[export]{adjustbox}
\usepackage{tabularx}
\usepackage{booktabs}

\usepackage{setspace}
\setstretch{1.2}

\makeatletter
\@ifundefined{KOMAClassName}{% if non-KOMA class
  \IfFileExists{parskip.sty}{%
    \usepackage{parskip}
  }{% else
    \setlength{\parindent}{0pt}
    \setlength{\parskip}{6pt plus 2pt minus 1pt}}
}{% if KOMA class
  \KOMAoptions{parskip=half}}
\makeatother
\usepackage{longtable,booktabs,array}
\newcounter{none} % for unnumbered tables
\usepackage{multirow}
\usepackage{calc} % for calculating minipage widths
% Correct order of tables after \paragraph or \subparagraph
\usepackage{etoolbox}
\makeatletter
\patchcmd\longtable{\par}{\if@noskipsec\mbox{}\fi\par}{}{}
\makeatother
% Allow footnotes in longtable head/foot
\IfFileExists{footnotehyper.sty}{\usepackage{footnotehyper}}{\usepackage{footnote}}
\makesavenoteenv{longtable}
\usepackage{graphicx}
\makeatletter
\newsavebox\pandoc@box
\newcommand*\pandocbounded[1]{% scales image to fit in text height/width
  \sbox\pandoc@box{#1}%
  \Gscale@div\@tempa{\textheight}{\dimexpr\ht\pandoc@box+\dp\pandoc@box\relax}%
  \Gscale@div\@tempb{\linewidth}{\wd\pandoc@box}%
  \ifdim\@tempb\p@<\@tempa\p@\let\@tempa\@tempb\fi% select the smaller of both
  \ifdim\@tempa\p@<\p@\scalebox{\@tempa}{\usebox\pandoc@box}%
  \else\usebox{\pandoc@box}%
  \fi%
}
% Set default figure placement to htbp
\def\fps@figure{htbp}
\makeatother
\setlength{\emergencystretch}{3em} % prevent overfull lines
\providecommand{\tightlist}{%
  \setlength{\itemsep}{0pt}\setlength{\parskip}{0pt}}


\hypersetup{colorlinks=true, linkcolor=blue, filecolor=magenta, urlcolor=cyan,}
\urlstyle{same}

\usepackage{colortbl}
\definecolor{table-row-color}{HTML}{999999}
\definecolor{table-rule-color}{HTML}{999999}
%\arrayrulecolor{black!40}
\arrayrulecolor{table-rule-color}     % color of \toprule, \midrule, \bottomrule
\setlength{\aboverulesep}{0pt}
\setlength{\belowrulesep}{0pt}


\setotherlanguages{english}
\IfFontExistsTF{Source Han Serif CN}
{\newfontfamily\chinesefont{Source Han Serif CN}}
{\IfFontExistsTF{Noto Serif CJK SC}
  {\newfontfamily\chinesefont{Noto Serif CJK SC}}
  {\IfFontExistsTF{SimSun}
    {\newfontfamily\chinesefont{SimSun}}
    {\IfFontExistsTF{FangSong}
      {\newfontfamily\chinesefont{FangSong}}
      {\newfontfamily\chinesefont{Arial Unicode MS}}
}}}
\IfFontExistsTF{Times New Roman}
{\newfontfamily\englishfont{Times New Roman}}
{\IfFontExistsTF{Liberation Serif}
  {\newfontfamily\englishfont{Liberation Serif}}
  {\IfFontExistsTF{DejaVu Serif}
    {\newfontfamily\englishfont{DejaVu Serif}}
    {\newfontfamily\englishfont{Arial}}
}}


\author{}
\date{}

\begin{document}


Article

\section{SF-YOLOv5: A Lightweight Small Object Detection AlgorithmBased
on Improved Feature Fusion
Mode}\label{sf-yolov5-a-lightweight-small-object-detection-algorithmbased-on-improved-feature-fusion-mode}

Haiying Liu 1,*, Fengqian Sun 1, Jason Gu 2 and Lixia Deng 1

1 School of Information and Automation Engineering, Qilu University of
Technology (Shandong Academyof Sciences), Jinan 250353, China

2 School of Electrical and Computer Engineering, Dalhousie University,
Halifax, NS B3H 4R2, Canada

\begin{itemize}
\tightlist
\item
  Correspondence: haiyingliu2019@qlu.edu.cn
\end{itemize}

Abstract: In the research of computer vision, a very challenging problem
is the detection of smallobjects. The existing detection algorithms
often focus on detecting full-scale objects, without makingproprietary
optimization for detecting small-size objects. For small objects dense
scenes, not onlythe accuracy is low, but also there is a certain waste
of computing resources. An improved detectionalgorithm was proposed for
small objects based on YOLOv5. By reasonably clipping the featuremap
output of the large object detection layer, the computing resources
required by the model weresignificantly reduced and the model becomes
more lightweight. An improved feature fusion method(PB-FPN) for small
object detection based on PANet and BiFPN was proposed, which
effectivelyincreased the detection ability for small object of the
algorithm. By introducing the spatial pyramidpooling (SPP) in the
backbone network into the feature fusion network and connecting with
themodel prediction head, the performance of the algorithm was
effectively enhanced. The experimentsdemonstrated that the improved
algorithm has very good results in detection accuracy and
real-timeability. Compared with the classical YOLOv5, the
\(\mathrm { m A P @ 0 . 5 }\) and
\(\operatorname* { m A P @ 0 . 5 : 0 . 9 5 }\) of SF-YOLOv5
wereincreased by \(1 . 6 \%\) and \(0 . 8 \%\) , respectively, the
number of parameters of the network were reduced by\(6 8 . 2 \%\) ,
computational resources (FLOPs) were reduced by \(1 2 . 7 \%\) , and the
inferring time of the modewas reduced by \(6 . 9 \%\) .

Keywords: object detection; visual tracking; small object; segmentation
and categorization; YOLO

\pandocbounded{\includegraphics[keepaspectratio,alt={image}]{images/a83da9e63d693963800fcf7ba55cf7290a97489aa171c4916da0b58a35ace0eb.jpg}}

Citation: Liu, H.; Sun, F.; Gu, J.; Deng,L. SF-YOLOv5: A Lightweight
SmallObject Detection Algorithm Based onImproved Feature Fusion
Mode.Sensors 2022, 22, 5817. https://doi.org/10.3390/s22155817

Academic Editor: Steven L.Waslander

Received: 11 July 2022

Accepted: 2 August 2022

Published: 4 August 2022

Publisher's Note: MDPI stays neutralwith regard to jurisdictional claims
inpublished maps and institutional affil-iations.

\pandocbounded{\includegraphics[keepaspectratio,alt={image}]{images/5317058b5cb7b648af44c23685169cc7f556633633d237c50bb8eab6843154fe.jpg}}

Copyright: © 2022 by the authors.Licensee MDPI, Basel, Switzerland.This
article is an open access articledistributed under the terms
andconditions of the Creative CommonsAttribution (CC BY) license
(https://creativecommons.org/licenses/by/4.0/).

\section{1. Introduction}\label{introduction}

Object detection aims to identify the location, size, spatial
relationship, and classesof specific targets in images or video
sequences. It is a cross domain of computer vision,digital image
processing, and machine vision. It always had high research and
applicationvalue in face recognition {[}1,2{]}, industrial defect
detection {[}3{]}, UAV aviation detection {[}4,5{]},traffic and vehicle
detection {[}6{]}, pedestrian detection and counting {[}7{]}, and other
fields.

In recent years, with the update and iteration of high-performance GPUs,
the con-tinuous emergence and improvement of various large-scale
datasets (such as DOTA {[}8{]},ImageNet {[}9{]}, COCO {[}10{]}) and the
strong rise of neural networks based on convolutional,object detection
algorithms using deep learning techniques have become the mainstreamof
contemporary object detection technology.

The current common object detection algorithms can be divided into
two-stage al-gorithm and one-stage detection algorithm according to
whether candidate regions aregenerated or not. The representative
algorithms of the former are: RCNN {[}11{]}, SPP(SpacePyramid
Pooling)-Net {[}12{]}, Fast-RCNN {[}13{]}, Faster-RCNN {[}14{]},
Mask-RCNN {[}15{]}, etc. Thelatter mainly includes YOLO (YOLOv1
{[}16{]}, YOLOv2 {[}17{]}, YOLOv3 {[}18{]}, YOLOv4 {[}19{]},YOLOv5
{[}20{]}, etc.) and SSD {[}21{]} algorithm. Although the two-stage
algorithms can reachhigh accuracy, the prolonged time of detection makes
it hard to satisfy the real-time require-ment in daily object detection
scenarios. Because it has the advantages of high precision

and fast speed, the one-stage detection algorithm has become the
research focus in thestudy of object detection.

For early versions of object detection algorithms, for example, YOLOv1
{[}16{]}, YOLOv2 {[}17{]},etc., its network structure is often the
stacking of multiple convolution layers and the full-connection layer.
Only the feature map of a fixed size is calculated, which greatly
limitedthe multi-scale detection ability of the model. On this occasion,
the detection accuracy of thealgorithm for small targets is not ideal.

In the case of solving the above problems, an effective method is to
adopt the imagepyramid strategy {[}22{]}, which can enhance the
multi-scale detection ability of the model;however, the high
computational cost and inference speed limit the development of theimage
pyramid strategy in an object detection algorithm.

Based on the aforementioned methods, FPN {[}23{]} up-samples the feature
map andthen fuses it with the corresponding output in the backbone
network in this process, andgenerates multiple new feature maps for
final detection. This modification improved thefeature expression
ability of the algorithm and further improved the effect of the
algorithmbased on real-time detection and multi-scale detection ability.
Meanwhile, FPN adoptedthe top-down feature fusion path, which also
strengthened the small object detectionperformance of the algorithm.

Due to the excellent performance of FPN, adding a feature fusion network
after thebackbone network of the model become the mainstream
configuration of object detectionalgorithms. For example, YOLOv3
{[}18{]}, YOLOv4 {[}19{]}, and YOLOv5 {[}20{]}, all use FPN as afeature
fusion network. The subsequent NAS-FPN {[}24{]}, ASFF {[}25{]}, PANet
{[}26{]}, BiFPN {[}27{]},and so on have achieved better results through
various feature fusion paths, but their basicideas are the same as that
of FPN.

Through the combination of well-designed backbones and feature fusion
networks,the current object detection algorithms can ensure the
detection accuracy and real-timecapability, and have the ability to
detect various scales objects; however, for the specific de-sign, these
algorithms often focus on various complex life scenarios. Their original
purposeis to guarantee that the mathematical model has perfect
generalization and detection abilityfor targets with different scales,
but ignore the detection optimization in some specificscale scenes. In
this context, designing a simple, efficient, lightweight, and
easy-to-deployalgorithm for small object detection has important study
significance.

YOLOv5 is the latest version of the YOLO series algorithms. Based on
YOLOv5, thisresearch has invented an improved YOLOv5 algorithm for small
object detection withthe name SF(Small-Fast)-YOLOv5, which can not only
significantly reduce the number ofparameters and calculated amount of
the model, but also has better performance comparedwith the original
version in the small object detection direction.

The major innovations and contributions of this research are as follows:

In the default network structure of YOLOv5, the network layer that
originally used togenerate the large object detection feature map has
been reasonably clipped. Whilerealizing the lightweight model, the
computing resources required by the model areeffectively released and
the speed of the model is greatly improved.

Based on the PANet {[}26{]} and BiFPN {[}27{]}, a new feature fusion
method (PB-FPN) forsmall object detection is proposed, which improved
the detection performance of thealgorithm for small targets.

The (SPP {[}12{]}) layer at the end of the backbone network is
introduced into the featurefusion network and connected with multiple
prediction heads to improve the featureexpression ability of the final
output feature map and further enhance the ability ofthe algorithm.

The other parts of this paper are arranged as follows: Section 2
introduces the ba-sic idea and network structure of the classical
algorithm YOLOv5; Section 3 illustratesthe improvement strategy of
SF-YOLOv5 in detail; Section 4 presents the experimentalenvironment,
dataset selection, experimental result analysis, ablation experiment,
algo-

rithm comparison, and algorithm generalization ability test; Finally,
Section 5 provides theconclusion, and introduces follow-up work and
improvement direction.

\section{2. Relevant Work}\label{relevant-work}

\section{2.1. Some Classical
Algorithms}\label{some-classical-algorithms}

The detection of small objects belongs to the problem of abnormal scale
in objectdetection. Most of the classical algorithms are full-size
detection algorithms, that is, thealgorithm itself will give priority to
ensuring the detection ability of objects of varioussizes, and make
various optimizations for small object detection on this basis.
Commonoptimization strategies include the using of image pyramids
{[}22{]} or FPN {[}23{]}. Addingmore detection heads {[}4{]} can also
attain a good effect, although it will greatly increase thecalculation
cost.

At present, common detection algorithms (such as {[}11--21{]}, etc.)
have been introduced.In addition, the latest YOLOv7 {[}28{]} (which is
still under update) and various improvedalgorithms based on ResNet
{[}29,30{]} also have excellent performance.

YOLOv5 is an algorithm with high reliability and stability, and it is
easy to deploy andtrain. At the same time, it is also one of the
one-stage detection algorithms with the highestaccuracy at present;
therefore, in this paper, we choose YOLOv5 for subsequent improve-ment.

\section{2.2. Basic Idea of YOLOv5}\label{basic-idea-of-yolov5}

YOLOv5 continues the consistent idea of the YOLO series in algorithm
design: namely,the image to be detected was processed through a input
layer (input) and sent to thebackbone for feature extraction. The
backbone obtains feature maps of different sizes, andthen fuses these
features through the feature fusion network (neck) to finally
generatethree feature maps P3, P4, and P5 (in the YOLOv5, the dimensions
are expressed withthe size of \(8 0 \times 8 0\) , \(4 0 \times 4 0\)
and \(2 0 \times 2 0\) ) to detect small, medium, and large objects
inthe picture, respectively. After the three feature maps were sent to
the prediction head(head), the confidence calculation and bounding-box
regression were executed for eachpixel in the feature map using the
preset prior anchor, so as to obtain a multi-dimensionalarray (BBoxes)
including object class, class confidence, box coordinates, width, and
heightinformation. By setting the corresponding thresholds
(confthreshold, objthreshold) tofilter the useless information in the
array, and performing a non-maximum suppression(NMS {[}31{]}) process,
the final detection information can be output. The process of
convertingthe input picture into BBoxes is called the inference process,
and the subsequent thresholdand NMS operations are called
post-processing. The post-processing does not involvethe network
structure. The default inference process of YOLOv5 can be represented
byFigure 1.

\pandocbounded{\includegraphics[keepaspectratio,alt={image}]{images/40dae9d57914a6acaf0649be2027477d5bc3c02799c5d7f58d34d912f1b354f7.jpg}}

Figure 1. The default inference flowchart of YOLOv5.

In the details of bounding-box regression, different from the previous
version of YOLOalgorithm, the process of YOLOv5 can be explained by (1).

\[
g _ {x} = 2 \sigma (s _ {x}) - 0. 5 + r _ {x}
\]

\[
\begin{array}{l} g _ {y} = 2 \sigma \left(s _ {y}\right) - 0. 5 + r _ {y} \\ g _ {h} = p _ {h} \left(2 \sigma \left(s _ {h}\right)\right) ^ {2} \end{array} \tag {1}
\]

\[
g _ {w} = p _ {w} \left(2 \sigma \left(s _ {w}\right)\right) ^ {2}
\]

In the above formula, set the coordinate value of the upper left corner
of the featuremap to (0, 0). \(r _ { x }\) and \(r _ { y }\) are the
unadjusted coordinates of the predicted center point.
\(g _ { x } ,\)\(g _ { y } , g _ { w } , g _ { h }\) represents the
information of the adjusted prediction box. \(p _ { w }\) and
\(p _ { h }\) are for theinformation of the prior anchor. \(s _ { x }\)
and \(s _ { y }\) represents the offset calculated by the model.The
process of adjusting the center coordinate and size of the preset prior
anchor to thecenter coordinate and size of the final prediction box is
called bounding-box regression.There are five versions of YOLOv5, namely
YOLOv5x, YOLOv5l, YOLOv5m, YOLOv5s,and YOLOv5n. The performance of each
version is shown in Table 1.

Table 1. Comprehensive performance of five versions of YOLOv5 on COCO
dataset {[}20{]}.

{\def\LTcaptype{none} % do not increment counter
\begin{longtable}[]{@{}|l|l|l|l|l|l|l|l|l|@{}}
\toprule\noalign{}
\endhead
\bottomrule\noalign{}
\endlastfoot
\hline
Model & Size (Pixels) & mAP @0.5:0.95 & mAP @0.5 & Time CPU b1 (ms) &
Time V100 b1 (ms) & Time V100 b32 (ms) & Params (M) & FLOPS @640 (B) \\
\hline
YOLOv5n & 640 & 28.0 & 45.7 & 45 & 6.6 & 0.6 & 1.9 & 4.5 \\
\hline
YOLOv5s & 640 & 37.4 & 56.8 & 98 & 6.4 & 0.9 & 7.2 & 16.5 \\
\hline
YOLOv5m & 640 & 45.4 & 64.1 & 224 & 8.2 & 1.7 & 21.2 & 49.0 \\
\hline
YOLOv5l & 640 & 49.0 & 67.3 & 430 & 10.1 & 2.7 & 46.5 & 109.1 \\
\hline
YOLOv5x & 640 & 50.7 & 68.9 & 766 & 12.1 & 4.8 & 86.7 & 205.7 \\
\hline
\end{longtable}
}

In this paper, for the comprehensive consideration of computing
resources, modelparameters, detection accuracy, deployment ability, and
algorithm practicability, we choseYOLOv5s as the basic algorithm and
made subsequent improvements and contributionsaround it.

\section{2.3. Network Structure of
YOLOv5}\label{network-structure-of-yolov5}

Generally speaking, the network structure of YOLOv5 refers to the
backbone and neck.

\section{2.3.1. Backbone}\label{backbone}

The backbone of YOLOv5 is shown in Figure 2. The main structure is the
stacking ofmultiple CBS (Conv \(^ +\) BatchNorm \(^ +\) SiLU) modules
and C3 modules, and finally one SPPFmodule is connected. CBS module is
used to assist C3 module in feature extraction, whileSPPF module
enhances the feature expression ability of the backbone.

\pandocbounded{\includegraphics[keepaspectratio,alt={image}]{images/edd0bbef14899b9ecce232145326b1b417daad4fd7fd89ac48d4693f168b0629.jpg}}

Figure 2. Default network structure of YOLOv5.

Therefore, in the backbone of YOLOv5, the most important layer is the C3
module.The basic idea of C3 comes from CSPNet (cross stage partial
networks {[}32{]}). C3 can actuallybe regarded as the specific
implementation of CSPNet. YOLOv5 uses the idea of CSPNetto build the C3
module, which not only ensures that the backbone has excellent
featureextraction ability, but also curbs the problem of gradient
information duplication in thebackbone.

\section{2.3.2. Neck}\label{neck}

In the neck, YOLOv5 uses the methods of FPN {[}23{]} and PAN {[}26{]},
as shown in Figure 3.The basic idea of FPN is to up-sampling the output
feature map (C3, C4, and C5) generatedby multiple convolution down
sampling operations from the feature extraction network togenerate
multiple new feature maps (P3, P4, and P5) for detecting different
scales targets.

\pandocbounded{\includegraphics[keepaspectratio,alt={image}]{images/9cc6afbe6776beedba79748a5d1768a95c65c139b6a00b175db1c2cd696550c2.jpg}}

Figure 3. The dotted line in the figure is the default feature fusion
path of YOLOv5.

The feature fusion path of FPN is top-down. On this basis, PAN
reintroduces a newbottom-up feature fusion path, which further enhance
the detection accuracy for differentscales objects.

\section{3. Approach}\label{approach}

The performance of YOLOv5 algorithm is very excellent for object
detection, but whenthe scenario is full of small targets, the detected
results cannot obtain ideal results, so thereis still has large room for
improvement. In order to solve these issues, some questionswere
proposed:

YOLOv5's feature extraction network set three different sizes of feature
map outputfor detection scenes of various scales. The process of
obtaining the feature maprequires multiple convolution down sampling
operations, which takes up a lot ofcomputing resources and parameters;
however, in the actual detection for smalltargets, is the down sampling
operation and feature fusion process for high-levelfeature map
necessary?

The feature fusion network of YOLOv5 combined a top-down path (FPN) with
abottom-up path (PAN), which enhances the detection performance on
various scaleobjects; however, for dense small object scenes, can we set
a new feature fusion pathto improve the feature expression potential of
the output feature map? Can we fusefeatures horizontally to further
enhance the detection performance of the algorithm?

YOLOv5 adds an SPPF module at the end of the backbone, which improves
theperformance of the feature extraction network through multiple
convolution andpooling operations with different sizes. Can we use this
idea to further tap the featureexpression potential of the feature map
in the end of the neck?

Aiming at the above three questions, we proposed a novel small object
detection algo-rithm: SF(Small-Fast)-YOLOv5. In the multiple down
sampling process of the backbone,in order to bring about the
lightweightness of the algorithm, we canceled the default C5(fifth
convolution down sampling operation) layer, and made corresponding
adjustmentsin the feature fusion part, leaving only the C3 and C4
feature maps for feature fusion. In thepart of the neck, based on PANet
and BiFPN, we redesigned a new small object detection

feature fusion path (PB-FPN), which aims to strengthen the fusion effect
of the model on theunderlying features and strengthen the small object
detection performance. Meanwhile, weintroduce the SPPF module that
originally existed only in the backbone at the connectionbetween the
neck and the head, which further tapped the feature expression potential
ofthe output feature maps.

\section{3.1. Feature Map Clipping}\label{feature-map-clipping}

YOLOv5 outputed three sizes of feature maps by default, which are used
to predictthe large, medium, and small objects, respectively. The way to
obtain these feature mapsis to carry out continuous convolution down
sampling operations and feature fusionprocesses. Table 2 shows the
numerical visualization results of YOLOv5 default featureextraction
network.

Table 2. Parameter of backbone in YOLOv5 network structure.

{\def\LTcaptype{none} % do not increment counter
\begin{longtable}[]{@{}|l|l|l|l|l|l|@{}}
\toprule\noalign{}
\endhead
\bottomrule\noalign{}
\endlastfoot
\hline
C & From & n & Params & Module & Arguments \\
\hline
0 & -1 & 1 & 3520 & CBS & {[}3, 32, 6, 2, 2{]} \\
\hline
1 & -1 & 1 & 18,560 & CBS & {[}32, 64, 3, 2{]} \\
\hline
2 & -1 & 1 & 18,816 & C3 & {[}64, 64, 1{]} \\
\hline
3 & -1 & 1 & 73,984 & CBS & {[}64, 128, 3, 2{]} \\
\hline
4 & -1 & 2 & 115,712 & C3 & {[}128, 128, 2{]} \\
\hline
5 & -1 & 1 & 295,424 & CBS & {[}128, 256, 3, 2{]} \\
\hline
6 & -1 & 3 & 625,152 & C3 & {[}256, 256, 3{]} \\
\hline
7 & -1 & 1 & 1,180,672 & CBS & {[}256, 512, 3, 2{]} \\
\hline
8 & -1 & 1 & 1,182,720 & C3 & {[}512, 512, 1{]} \\
\hline
9 & -1 & 1 & 656,896 & SPPF & {[}512, 512, 5{]} \\
\hline
\end{longtable}
}

As shown in Table 2, C represents the position of the module, From
\(^ { - 1 }\) indicates thatthe input of this layer comes from the
previous layer, n indicates the amount of modules,params indicates the
parameter usage of this layer, module indicates the name of the mod-ule,
arguments represent the input channel value, output channel value, and
convolutionkernel attribute, respectively. The 7 and 8 layers indicates
the C5 sampling layer in the fea-ture extraction network; as can be
seen, the sampling layer occupies enormous parameters.

In the proposed mathematical model, for the purpose of being
lightweight, the defaultC5 feature map of the traditional YOLOv5 is
deleted. The feature map clipping corre-sponding to the improved method
in this section can be represented by the flowchart inFigure 4.

\pandocbounded{\includegraphics[keepaspectratio,alt={image}]{images/1177682839fec199f1fbd351573c182d8a51e75779d39e15636e664ef28856ed.jpg}}

Figure 4. The inference flowchart of SF-YOLOv5. The backbone has been
cut accordingly, correspond-ing to the improvement method proposed in
Section 3.1, neck (PBS) represents a new feature fusionnetwork with
PB-FPN and SPPF prediction head, which corresponds to the methods
described inSections 3.2 and 3.3 of the improved algorithm in this
paper.

\section{3.2. Improvement for Feature Fusion Path
(PB-FPN)}\label{improvement-for-feature-fusion-path-pb-fpn}

Early object detection algorithms usually used the high-level feature
map that gen-erated by multiple down sampling operations for detection,
which will make the modelunable to effectively predict objects of
various scales. The introduction of FPN {[}23{]} solvesthis problem
well, and its basic process is shown in Figure 5b.

\pandocbounded{\includegraphics[keepaspectratio,alt={image}]{images/7e857d8989ccbfc2c185fae9f6357849eb7afa4cd0b718b54a270b93f4fc9f47.jpg}}

\begin{enumerate}
\def\labelenumi{(\alph{enumi})}
\tightlist
\item
  Default
\end{enumerate}

\pandocbounded{\includegraphics[keepaspectratio,alt={image}]{images/74beaa7866122a96878ceea3971ab77d963d85853f7d989085ec3c339439f27d.jpg}}

\begin{enumerate}
\def\labelenumi{(\alph{enumi})}
\setcounter{enumi}{1}
\tightlist
\item
  FPN
\end{enumerate}

\pandocbounded{\includegraphics[keepaspectratio,alt={image}]{images/9cefd1d12a50bd483943fc5456d54fda6170a3ee52c806a183f7db528288758a.jpg}}

\begin{enumerate}
\def\labelenumi{(\alph{enumi})}
\setcounter{enumi}{2}
\tightlist
\item
  PAN
\end{enumerate}

\pandocbounded{\includegraphics[keepaspectratio,alt={image}]{images/93192fd31bfd7665d2cfe18f47537fffd11bee931aa5d3740a52c5bee5d92783.jpg}}

\begin{enumerate}
\def\labelenumi{(\alph{enumi})}
\setcounter{enumi}{3}
\tightlist
\item
  BiFPN
\end{enumerate}

Figure 5. Some common feature fusion paths.

Taking the P4 node in Figure 5d as an example:

\[
P _ {4} ^ {t d} = \operatorname {C o n v} \left(\frac {w _ {1} \cdot P _ {4} ^ {i n} + w _ {2} \cdot \operatorname {R e s i z e} \left(P _ {5} ^ {i n}\right)}{w _ {1} + w _ {2} + \varepsilon}\right) \tag {2}
\]

\[
P _ {4} ^ {\text {o u t}} = \operatorname {C o n v} \left(\frac {w _ {1} ^ {\prime} \cdot P _ {4} ^ {\text {i n}} + w _ {2} ^ {\prime} \cdot P _ {4} ^ {\text {t d}} + w _ {3} ^ {\prime} \cdot \operatorname {R e s i z e} \left(P _ {3} ^ {\text {o u t}}\right)}{w _ {1} ^ {\prime} + w _ {2} ^ {\prime} + w _ {3} ^ {\prime} + \varepsilon}\right) \tag {3}
\]

In the (2) and (3), \(P _ { 4 } ^ { t d }\) represents the intermediate
feature map between C4 and P4, \(P _ { 4 } ^ { i n }\)and
\(P _ { 4 } ^ { o u t }\) 4 4correspond to C4 and P4 respectively, Conv
and Resize correspond to convolutionand sampling operation respectively,
\(w\) and ε represent weight and a preset small value toavoid numerical
instability, which is usually set to 0.0001.

Based on the previous analysis, we designed an improved feature fusion
path, asshown in Figure 6:

\pandocbounded{\includegraphics[keepaspectratio,alt={image}]{images/8f2fd7e257f84657057c59544141b814881cdb3e70a812064c83d4ba12e19577.jpg}}

\begin{enumerate}
\def\labelenumi{(\alph{enumi})}
\tightlist
\item
  Default PB-FPN
\end{enumerate}

\pandocbounded{\includegraphics[keepaspectratio,alt={image}]{images/ae5f4336855747e45e3bb76aac59df159e5f36e742269a347188e8af6970ac96.jpg}}

\begin{enumerate}
\def\labelenumi{(\alph{enumi})}
\setcounter{enumi}{1}
\tightlist
\item
  PB-FPN after cutting C5
\end{enumerate}

Figure 6. The feature fusion path of PB-FPN.

Compared with the classical algorithms PANet {[}26{]} and BiFPN
{[}27{]}, this sectionadditionally introduces a new feature fusion path
to further integrate the features from thehigh-level into the bottom.
Meanwhile, new branches are set horizontally to participatein the fusion
process of the bottom, which further enhances the detection effect of
thealgorithm for small targets.

\section{3.3. The Improved Feature Fusion Network
(SPPF)}\label{the-improved-feature-fusion-network-sppf}

The latest version of YOLOv5 used the SPPF (SPP-Fast) module. Compared
with theoriginal SPP module, its structure comparison is shown in Figure
7. By simplifying thepooling process, SPPF avoided the repeated
operation of SPP and significantly improvedthe running speed of the
module.

\pandocbounded{\includegraphics[keepaspectratio,alt={image}]{images/768359bc742a3969c7bd8a3a0456147c72932ef963f96ec52a489faca62dfa86.jpg}}

\begin{enumerate}
\def\labelenumi{(\alph{enumi})}
\tightlist
\item
  SPPF
\end{enumerate}

\pandocbounded{\includegraphics[keepaspectratio,alt={image}]{images/847dc863a52810440da66cb23880073450f26e9f8aff4dd7ed1c9d430c06b99f.jpg}}

\begin{enumerate}
\def\labelenumi{(\alph{enumi})}
\setcounter{enumi}{1}
\tightlist
\item
  SPP
\end{enumerate}

Figure 7. Structure comparison of SPPF and SPP.

In this paper, we introduced several SPPF modules at the connection
between thefeature fusion network and the model prediction head, in
order to further tap the featureexpression potential of the finally
output feature map by the neck and sent to the head, andfurther enhance
the performance of the model. The modified network in this section can
beexpressed in Figure 8.

\pandocbounded{\includegraphics[keepaspectratio,alt={image}]{images/ac178bce89ae4050bdf97d5d4ef0458f8960bfbdbdb961950e3a898815c66b94.jpg}}

Figure 8. Structure diagram after adding SPPF at the junction of feature
fusion network and predic-tion head.

\section{4. Experiments}\label{experiments}

\section{4.1. Experimental Environment}\label{experimental-environment}

In this experiment, the workstation is win10, 16 GB, the CPU is used
AMD-4800H,the GPU is used Nvidia-Rtx2060. The algorithm is based on
PyTorch, using CUDA11.1 tooperation acceleration. The training epoch is
set to 300, and the mosaic data enhancementand a prior anchor adaptive
adjustment strategy are used.

\section{4.2. Dataset Introduction}\label{dataset-introduction}

In order to testify the performance of the improved algorithm
(SF-YOLOv5) for thesmall object detection, this study adopted the WIDER
FACE {[}33{]} dataset to train and verifythe algorithm. This dataset
contains the annotation information of 393,703 faces. As shownin Figure
9, the dataset is highly variable in scale, posture, angle, light, and
occlusion. Thedataset is complex and includes a great quantity of dense
small targets.

Scale

\pandocbounded{\includegraphics[keepaspectratio,alt={image}]{images/2e732552f81187b63e07b1d2cefd51df53bfcc9b5ca5fc9b6e4d45c7402e49fb.jpg}}

Pose

\pandocbounded{\includegraphics[keepaspectratio,alt={image}]{images/dbba54665d8ac8e92e91fb967cbde60e6e2b0cd4f18c6eff29b9f58122d8613b.jpg}}

Occlusion

\pandocbounded{\includegraphics[keepaspectratio,alt={image}]{images/07fd662c0e120a43232ebf100f94c415e70893653b6964905dd19b1250795fe1.jpg}}

Expression

\pandocbounded{\includegraphics[keepaspectratio,alt={image}]{images/571d0881c55e4eb9a09c9b710902fa13a0de74c707ea6ff021987c9979afdb57.jpg}}

Makeup

\pandocbounded{\includegraphics[keepaspectratio,alt={image}]{images/b99a58791abe899c605fe26a564073b97d3cdd070905f955691d783e08c530a4.jpg}}

Illumination

\pandocbounded{\includegraphics[keepaspectratio,alt={image}]{images/d40beb48aca735890494062461412c549b42e6c1baaafe5a6b7c07ad21539a12.jpg}}

Figure 9. The WIDER FACE dataset contains a variety of complex and rich
scenes {[}33{]}.

Considering that the SF-YOLOv5 mainly detects small objects, while the
originalWIDER FACE dataset contains not only small targets, but also
large targets, which willaffect the experimental results; therefore, we
first screened according to the number andsize of face images in a
single image in the original dataset. In total, 4441 images were usedto
train the algorithm and 1123 images were used to verify the algorithm.
The numberproportion of the train and verify was 4:1, same as the
original dataset. Figure 10 is anattribute visualization result of the
new dataset.

\pandocbounded{\includegraphics[keepaspectratio,alt={image}]{images/261150ae917a4dca95ec15909ae3207951c521fed35283f02427d2c842d4ac36.jpg}}

(a)the number of dataset labels

\pandocbounded{\includegraphics[keepaspectratio,alt={image}]{images/f25ffbef3f809614371ab8d16910f71fd4af190ef1596361fccdd0badd7c823a.jpg}}

(b)the label location of dataset

\pandocbounded{\includegraphics[keepaspectratio,alt={image}]{images/f450d1804abf3ced12c55d5d4e923a60ba08e0f959786d78def718f67ebdb652.jpg}}

(c)the label size of dataset

Figure 10. The attribute visualization results of the dataset used in
this paper.

Figure 10a expressed the number of labels in the dataset. Figure 10b
shows thecentral coordinate position of the object. Figure 10c shows the
size of the object. It can beseen that this dataset is sufficient to
ensure the training and verification of small objectdetection algorithm.

\section{4.3. Evaluation Criterion}\label{evaluation-criterion}

At present, the mainstream general indicators for evaluating the
performance of objectdetection algorithms include precision, recall, AP
(average precision), mAP (mean AP),parameters(the number of parameters
in model), FLOPs (floating point operations persecond), inference time,
etc.

Considering that mAP can more reasonably reflect the accuracy of the
algorithm thanprecision and recall, in this experiment, we choose mAP
\(@ 0 . 5\) , mAP@0.5:0.95, parameters,FLOPs, and inference time to
evaluate the algorithm.

\section{4.4. Analysis of Experimental
Results}\label{analysis-of-experimental-results}

The experimental performance comparison between YOLOv5s and SF-YOLOv5
isillustrated in Table 3. Compared with the traditional algorithm
YOLOv5s, the mAP@0.5 andmAP@0.5:0.95 of SF-YOLOv5 has been increased by
1.6 and 0.8, respectively, which provedthe improvement in comprehensive
detection performance of SF-YOLOv5. At the sametime, the parameters (M)
value and FLOPs (G) value of SF-YOLOv5 are reduced by \(6 8 . 2 \%\)and
\(1 2 . 7 \%\) , respectively, indicating that SF-YOLOv5 can further
decrease the number ofparameters and lessen computing power required for
model operation under the conditionof improving the detection accuracy,
which enhances the deployment performance of thealgorithm on various
low-end workstations and small mobile devices. The reduction ininference
time (ms) shows that the model can process more images and videos in the
sametime, which ensures the speed of SF-YOLOv5 in detecting small
objects.

\section{4.5. Ablation Experiment}\label{ablation-experiment}

Table 3 expressed the ablation experimental data of improved methods 1,
2, and 3proposed in this paper. YOLOv5-N5 is a new algorithm obtained by
cutting the featurelayer with the default YOLOv5, which corresponds to
the improved method 1 proposed inthis paper. It is not difficult to find
that although this improvement has increased by 0.2in mAP@0.5, the value
of mAP@0.5:0.95 has decreased by 0.2. The calculation formula ofGIOU
{[}34{]} used by YOLOv5 is shown in (4).

\[
L = 1 - I o U + \frac {\left| A - P \cup P ^ {g t} \right|}{\left| A \right|} \tag {4}
\]

where \(P\) and \(P ^ { g t }\) express the area of the prediction box
and ground truth box, respectively,and \(A\) is the minimum area of the
area composed of \(P\) and \(P ^ { g t }\) . We hope the value of \(L\)
tobe infinitely close to 0. In this case, the lower the value of IOU,
the lower the requirementfor the position coincidence when verifying the
accuracy of the algorithm.

Table 3. Performance comparison of SF-YOLOv5 and other algorithms on
small object datasets.

{\def\LTcaptype{none} % do not increment counter
\begin{longtable}[]{@{}|l|l|l|l|l|l|l|@{}}
\toprule\noalign{}
\endhead
\bottomrule\noalign{}
\endlastfoot
\hline
Methods & Size & mAP@0.5 & mAP@0.5:0.95 & Parameters (M) & FLOPs (G) &
Inference Time (ms) \\
\hline
YOLOv5s & 640 & 69.7 & 35.5 & 7.01 & 15.8 & 13.1 \\
\hline
YOLOv5-N5 & 640 & 69.9 & 35.3 & 1.88 & 11.2 & 10.4 \\
\hline
YOLOv5-PB & 640 & 70.8 & 36.0 & 2.03 & 12.7 & 11.3 \\
\hline
SF-YOLOv5 & 640 & 71.3 & 36.3 & 2.23 & 13.8 & 12.2 \\
\hline
improvement & - & +1.6 & +0.8 & -68.2\% & -12.7\% & -6.9\% \\
\hline
YOLOv5n & 640 & 63.9 & 30.8 & 1.76 & 4.20 & 8.60 \\
\hline
SF-YOLOv5 & 640 & 71.3 & 36.3 & 2.23 & 13.8 & 12.2 \\
\hline
improvement & - & +7.4 & +5.5 & +26.7\% & +228.6\% & +41.9\% \\
\hline
YOLOv7-tiny & 640 & 68.4 & 33.0 & 6.01 & 13.0 & 13.9 \\
\hline
SF-YOLOv5 & 640 & 71.3 & 36.3 & 2.23 & 13.8 & 12.2 \\
\hline
improvement & - & +2.9 & +3.3 & -62.9\% & +6.2\% & -12.2\% \\
\hline
YOLOv3 & 640 & 74.5 & 39.5 & 61.5 & 154.7 & 44.7 \\
\hline
SF-YOLOv5L & 640 & 75.1 & 39.7 & 15.5 & 91.2 & 49.4 \\
\hline
improvement & - & +0.6 & +0.2 & -74.8\% & -41.0\% & +10.5\% \\
\hline
YOLOv7 & 640 & 76.1 & 39.5 & 36.5 & 103.2 & 18.0 \\
\hline
SF-YOLOv5L & 640 & 75.1 & 39.7 & 15.5 & 91.2 & 49.4 \\
\hline
improvement & - & -1.0 & +0.2 & -57.5\% & -11.6\% & +174.4\% \\
\hline
ResNeXt-CSP & 640 & 73.7 & 37.6 & 31.8 & 58.9 & 32.6 \\
\hline
SF-YOLOv5L & 640 & 75.1 & 39.7 & 15.5 & 91.2 & 49.4 \\
\hline
improvement & - & +1.4 & +2.1 & -51.3\% & +54.8\% & +51.5\% \\
\hline
\end{longtable}
}

It can be seen from Table 3 and formula (4), the improved method 1 can
achievesignificant weight improvements after removing the high-level
feature map. Although thedetection performance has been slightly
affected: the model can identify more objects; theposition of the
prediction box will deviate to a certain extent.

YOLOv5-PB is a new algorithm obtained by setting the feature fusion path
as PB-FPNon the basis of the improved method 1, corresponding to the
improved method 2 in thispaper. Compared with YOLOv5-N5, this algorithm
slightly increases the computationalresources and number of parameters,
but greatly enhances the detection accuracy, the valueof
\(\mathrm { m A P @ 0 . 5 }\) and mAP@0.5:0.95 by 0.9 and 0.7,
respectively. It proves the feasibility ofintroducing a more efficient
feature fusion path to enhance the ability of algorithm.

The final SF-YOLOV5 was obtained by adding SPPF module at the junction
betweenthe neck and head of YOLOV5-PB, corresponding to the improved
method 3. Comparedwith the improved method 2, YOLOV5-PB slightly
increases the number of parameters andfurther enhances the detection
accuracy of the algorithm.

Figure 11 shows the visualization of ablation experimental data of
improved methods1, 2, and 3. Figure 12 shows the detection results of
SF-YOLOv5.

\pandocbounded{\includegraphics[keepaspectratio,alt={image}]{images/8c5ced4ff494869817e2c76a1f5a10fc0c9ca49f75168b8f242652c81ddce725.jpg}}

(a)mAP@0.5

\pandocbounded{\includegraphics[keepaspectratio,alt={image}]{images/c52bff19510d84a0f0a952668a7c727921368fe59bce9f814763fbd76462e2b0.jpg}}

(b)mAP@0.5:0.95

\pandocbounded{\includegraphics[keepaspectratio,alt={image}]{images/5ebb02bccce762633deade1daca8256f074e440dddc943510592d29006339b16.jpg}}

(c)Loss

Figure 11. The ablation experimental results of the improved methods 1,
2, and 3.

\pandocbounded{\includegraphics[keepaspectratio,alt={image}]{images/a6c470c56c016b5fa9f2e9347799d04d13ef5058488f477d31c0e668d4f6b0b8.jpg}}

Figure 12. Detection effect of SF-YOLOv5 on WIDER FACE dataset.

\section{4.6. Comparison with Other Classical
Algorithms}\label{comparison-with-other-classical-algorithms}

We compare SF-YOLOv5 with several classical detection algorithms, and
the resultsare expressed in Table 3. Among them, YOLOv5n is the latest
lightweight algorithm ofYOLOv5, and YOLOv3 is a relatively mature
large-scale one-stage detection algorithm.YOLOv7 {[}28{]} is the latest
algorithm of YOLO family at present, which has the
strongestcomprehensive performance in full-scale detection, and
YOLOv7-tiny is a lightweightversion of YOLOv7, which has similar
parameter quantities and calculation quantitieswith SF-YOLOv5.
ResNeXt-CSP is a new detector combined with classical algorithmsResNeXt
{[}29{]} and CSPNet {[}32{]}, which have excellent performance. It is
not difficult to findthat although YOLOv5n is lighter and faster, the
detection accuracy is too low comparedwith SF-YOLOv5. Compared with
YOLOv7-tiny, SF-YOLOv5 has advantages in detectionaccuracy, speed, and
lightweight effect in the small object dataset, but it slightly
increasesthe amount of calculation.

Because YOLOv3, YOLOv7, and ResNeXt-CSP are not lightweight algorithms,
it isdifficult to directly compare with SF-YLOv5, so we adjusted the
scaling coefficient of SF-YOLOv5. This operation is adopted in YOLOv4,
YOLOv5, and YOLOv7. The default value(0.33, 0.25) is adjusted to the
same (1, 1) as YOLOv7. The adjusted SF-YOLOv5 was namedSF-YOLOv5L. It
has the same network structure and improvement idea as SF-YOLOv5,and the
only difference is in the model size. The comparison of these algorithms
is shownin Table 3. On the whole, SF-YOLOv5L is better than YOLOv3,
ResNeXt-CSP, and itsperformance is close to that of the latest YOLOv7.
This proved the feasibility and reliabilityof the improved method
proposed in this paper.

\section{4.7. Performance of Novel SF-YOLOv5 on Other
Datasets}\label{performance-of-novel-sf-yolov5-on-other-datasets}

In this section, we further prove the ability of SF-YOLOv5 on other
datasets, so asto judge whether the improvement made by SF-YOLOv5 in the
direction of small objectdetection is universal.

TinyPerson {[}35{]}, VisDrone {[}36{]}, and VOC2012 {[}37,38{]} public
datasets were selectedfor testing experiment in this section. Among
them, TinyPerson and VisDrone contain alarge number of small targets,
which is suitable for the verification of SF-YOLOv5 proposedin this
paper. Because there are not only small targets, but also a large number
of largetargets in VOC2012, it is mainly used to verify the performance
of improved methods 2and 3 proposed in this paper in the full-scale
detection; therefore, in the experiment forVOC2012, we fine tuned the
prediction head of SF-YOLOv5 and re added the P5 detectionhead. The
algorithm is named SF-YOLOv5-P5. Table 4 shows the performance
comparisonof SF-YOLOv5 and YOLOv5s on these datasets.

Table 4. Performance comparison between YOLOv5s and SF-YOLOv5 on other
datasets.

{\def\LTcaptype{none} % do not increment counter
\begin{longtable}[]{@{}llllllll@{}}
\toprule\noalign{}
\endhead
\bottomrule\noalign{}
\endlastfoot
Dataset & Methods & Size & mAP@0.5 & mAP@0.5:0.95 & Parameters (M) &
FLOPs (G) & Inference Time (ms) \\
\multirow{3}{*}{TinyPerson} & YOLOv5s & 640 & 18.7 & 6.0 & 7.02 & 15.8 &
12.8 \\
& SF-YOLOv5 & 640 & 20.0 & 6.5 & 2.23 & 13.8 & 11.0 \\
& improvement & - & +1.3 & +0.5 & -68.2\% & -12.7\% & -14.1\% \\
\multirow{3}{*}{VisDrone} & YOLOv5s & 640 & 33.0 & 17.9 & 7.04 & 15.9 &
12.6 \\
& SF-YOLOv5 & 640 & 34.3 & 18.2 & 2.24 & 13.8 & 11.5 \\
& improvement & - & +1.3 & +0.3 & -68.2\% & -13.2\% & -8.7\% \\
\multirow{3}{*}{VOC2012} & YOLOv5s & 640 & 60.8 & 37.0 & 7.06 & 15.9 &
25.8 \\
& SF-YOLOv5-P5 & 640 & 61.2 & 38.3 & 4.59 & 15.7 & 24.8 \\
& improvement & - & +0.4 & +1.3 & -34.9\% & -1.3\% & -3.9\% \\
\end{longtable}
}

To sum up, SF-YOLOv5 still achieved improvement in detection accuracy
and compre-hensive performance on TinyPerson, VisDrone and VOC2012
datasets. It is proved that theimproved methods 1, 2, and 3 have obvious
improvement effect in the direction for smallobject detection and have
certain versatility in full-scale detection. Other experimentalresults
in this section are shown in Figures 13 and 14.

\pandocbounded{\includegraphics[keepaspectratio,alt={image}]{images/cb452bc277c9ffcc8341ffc046d271856557f691d7c5fca2d03a3304a06c3c32.jpg}}

\begin{enumerate}
\def\labelenumi{(\alph{enumi})}
\tightlist
\item
  TinyPerson
\end{enumerate}

\pandocbounded{\includegraphics[keepaspectratio,alt={image}]{images/9a18bb5ad1ed477c0c368bfab1fcaa274a25686a783efb15fe61973afa931f72.jpg}}

\begin{enumerate}
\def\labelenumi{(\alph{enumi})}
\setcounter{enumi}{1}
\tightlist
\item
  VisDrone
\end{enumerate}

\pandocbounded{\includegraphics[keepaspectratio,alt={image}]{images/0210c8fc55e9926cd9ce82519cdb3a1b0df653e8e9dcdc871c119417fa190caa.jpg}}

\begin{enumerate}
\def\labelenumi{(\alph{enumi})}
\setcounter{enumi}{2}
\tightlist
\item
  VOC2012
\end{enumerate}

Figure 13. Comparison of detection accuracy between SF-YOLOv5 and
YOLOv5s on other datasets.

\pandocbounded{\includegraphics[keepaspectratio,alt={image}]{images/2b2d2dff51632b56456c0aa4c008d98b3b2b8115c04fce13684c999c42277772.jpg}}

\pandocbounded{\includegraphics[keepaspectratio,alt={image}]{images/5c795e6bc7b73a9fba0479fbf00f11e7fa76f29eb44ffca6d931d4284ac6bf45.jpg}}

\pandocbounded{\includegraphics[keepaspectratio,alt={image}]{images/38005e8b39650441db71674d84c10202081715022b815ca63300820e47a83e29.jpg}}

Figure 14. Detection effect of SF-YOLOv5 on other dataset.

\section{5. Conclusions}\label{conclusions}

In this paper, a lightweight small object detection algorithm based on
an improvedfeature fusion mode is proposed, which is dedicated to
improving the detection effect onsmall targets. When the algorithm is
significantly lighter, the detection accuracy is alsosignificantly
improved. By reducing the convolution down sampling operation in
thenetwork structure and trimming the output feature map, the model
parameter amountand computation amount are significantly reduced under
the premise of ensuring thedetection accuracy. Ablation experiments show
that the contribution of high-level featuremaps generated by multiple
convolution downsampling can be replaced by more effectivefeature fusion
methods when detecting small objects. Besides this, in order to
improvethe small object detection accuracy of the model, a new feature
fusion method (PB-FPN)is proposed based on PANet and BiFPN. By setting a
fast SPP operation at the junctionof the neck and the head, the feature
expression potential of the output feature map ofthe backbone is fully
exploited, and the detection performance of the algorithm is further

improved. The revised model has good generalization ability and can also
be applied tosmall object detection scenarios of other types of targets.

In the follow-up research, we will continue to improve the detection
effect of thealgorithm on ultra-small targets, and we will further
explore the detection potential ofthe small object detection feature
fusion method (PB-FPN) proposed in this paper onfull-scale targets.

Author Contributions: Supervision, H.L.; Writing---original draft, F.S.;
Writing---review and editing,J.G.; Validation, L.D. All authors have
read and agreed to the published version of the manuscript.

Funding: This research was funded by QLUTGJHZ2018019.

Conflicts of Interest: The authors declare no conflict of interest.

\section{References}\label{references}

\begin{enumerate}
\def\labelenumi{\arabic{enumi}.}
\item
  Kowalski, M.; Grudzie ´n, A.; Mierzejewski, K. Thermal--Visible Face
  Recognition Based on CNN Features and Triple TripletConfiguration for
  On-the-Move Identity Verification. Sensors 2022, 22, 5012.
  {[}CrossRef{]} {[}PubMed{]}
\item
  Zhang, Z.; Shen, W.; Qiao, S.; Wang, Y.; Wang, B.; Yuille, A. Robust
  face detection via learning small faces on hard images. InProceedings
  of the IEEE/CVF Winter Conference on Applications of Computer Vision
  (WACV), Snowmass Village, CO, USA,1--5 March 2020; pp.~1361--1370.
\item
  Bai, T.; Gao, J.; Yang, J.; Yao, D. A study on railway surface defects
  detection based on machine vision. Entropy 2021, 23,
  1437.{[}CrossRef{]} {[}PubMed{]}
\item
  Zhu, X.; Lyu, S.; Wang, X.; Zhao, Q. TPH-YOLOv5: Improved YOLOv5 based
  on transformer prediction head for object detectionon drone-captured
  scenarios. In Proceedings of the IEEE/CVF International Conference on
  Computer Vision (ICCV), Montreal,BC, Canada, 11--17 October 2021;
  pp.~2778--2788.
\item
  Zhang, P.; Zhong, Y.; Li, X. SlimYOLOv3: Narrower, faster and better
  for real-time UAV applications. In Proceedings of theIEEE/CVF
  International Conference on Computer Vision (ICCV), Seoul, Korea,
  27--28 October 2019.
\item
  Gu, Y.; Si, B. A novel lightweight real-time traffic sign detection
  integration framework based on YOLOv4. Entropy 2022, 24,
  487.{[}CrossRef{]} {[}PubMed{]}
\item
  Liu, W.; Liao, S.; Ren, W.; Hu, W.; Yu, Y. High-level semantic feature
  detection: A new perspective for pedestrian detection.In Proceedings
  of the IEEE/CVF Conference on Computer Vision and Pattern Recognition
  (CVPR), Long Beach, CA, USA,16--20 June 2019; pp.~5187--5196.
\item
  Xia, G.S.; Bai, X.; Ding, J.; Zhu, Z.; Belongie, S.; Luo, J.; Datcu,
  M.; Pelillo, M.; Zhang, L. DOTA: A large-scale dataset for
  objectdetection in aerial images. In Proceedings of the IEEE
  Conference on Computer Vision and Pattern Recognition (CVPR), SaltLake
  City, UT, USA, 18--23 June 2018; pp.~3974--3983.
\item
  Deng, J.; Dong, W.; Socher, R.; Li, L.J.; Li, K.; Fei-Fei, L.
  Imagenet: A large-scale hierarchical image database. In Proceedings
  ofthe IEEE Conference on Computer Vision and Pattern Recognition
  (CVPR), Miami, FL, USA, 20--25 June 2009; pp.~248--255.
\item
  Lin, T.Y.; Maire, M.; Belongie, S.; Hays, J.; Perona, P.; Ramanan, D.;
  Dollár, P.; Zitnick, C.L. Microsoft coco: Common objects incontext. In
  Proceedings of the European Conference on Computer Vision (ECCV),
  Zurich, Switzerland, 6--12 September 2014;pp.~740--755.
\item
  Girshick, R.; Donahue, J.; Darrell, T.; Malik, J. Rich feature
  hierarchies for accurate object detection and semantic segmentation.In
  Proceedings of the IEEE Conference on Computer Vision and Pattern
  Recognition (CVPR), Columbus, OH, USA, 23--28 June2014; pp.~580--587.
\item
  He, K.; Zhang, X.; Ren, S.; Sun, J. Spatial pyramid pooling in deep
  convolutional networks for visual recognition. IEEE Trans.Pattern
  Anal. Mach. Intell. 2015, 37, 1904--1916. {[}CrossRef{]} {[}PubMed{]}
\item
  Girshick, R. Fast R-CNN. In Proceedings of the IEEE International
  Conference on Computer Vision (ICCV), Santiago, Chile,7--13 December
  2015; pp.~1440--1448.
\item
  Ren, S.; He, K.; Girshick, R.; Sun, J. Faster r-cnn: Towards real-time
  object detection with region proposal networks. Adv. NeuralInf.
  Process. Syst. 2015, 28, 91--99. {[}CrossRef{]} {[}PubMed{]}
\item
  He, K.; Gkioxari, G.; Dollár, P.; Girshick, R. Mask r-cnn. In
  Proceedings of the IEEE International Conference on Computer
  Vision(ICCV), Venice, Italy, 22--29 October 2017; pp.~2961--2969.
\item
  Redmon, J.; Divvala, S.; Girshick, R.; Farhadi, A. You only look once:
  Unified, real-time object detection. In Proceedings of theIEEE
  Conference on Computer Vision and Pattern Recognition (CVPR), Las
  Vegas, NV, USA, 27--30 June 2016; pp.~779--788.
\item
  Redmon, J.; Farhadi, A. YOLO9000: Better, faster, stronger. In
  Proceedings of the IEEE Conference on Computer Vision andPattern
  Recognition (CVPR), Honolulu, HI, USA, 21--26 July 2017;
  pp.~7263--7271.
\item
  Redmon, J.; Farhadi, A. Yolov3: An incremental improvement. arXiv
  2018, arXiv:1804.02767.
\item
  Bochkovskiy, A.; Wang, C.Y.; Liao, H.Y.M. Yolov4: Optimal speed and
  accuracy of object detection. arXiv 2020, arXiv:2004.10934.
\item
  Glenn, J. Yolov5-6.1---TensorRT, TensorFlow Edge TPU and OpenVINO
  Export and Inference. 2022. Available online:
  https://github.com/ultralytics/yolov5/releases/tag/v6.1 (accessed on
  22 February 2022).
\item
  Liu, W.; Anguelov, D.; Erhan, D.; Szegedy, C.; Reed, S.; Fu, C.Y.;
  Berg, A.C. Ssd: Single shot multibox detector. In Proceedings ofthe
  European conference on computer vision (ECCV), Amsterdam, The
  Netherlands, 11--14 October 2016; pp.~21--37.
\item
  Singh, B.; Davis, L.S. An analysis of scale invariance in object
  detection snip. In Proceedings of the IEEE Conference on
  ComputerVision and Pattern Recognition (CVPR), Salt Lake City, UT,
  USA, 18--23 June 2018; pp.~3578--3587.
\item
  Lin, T.Y.; Dollár, P.; Girshick, R.; He, K.; Hariharan, B.; Belongie,
  S. Feature pyramid networks for object detection. In Proceedingsof the
  IEEE Conference on Computer Vision and Pattern Recognition (CVPR),
  Honolulu, HI, USA, 21--26 July 2017; pp.~2117--2125.
\item
  Ghiasi, G.; Lin, T.Y.; Le, Q.V. Nas-fpn: Learning scalable feature
  pyramid architecture for object detection. In Proceedingsof the
  IEEE/CVF Conference on Computer Vision and Pattern Recognition (CVPR),
  Long Beach, CA, USA, 16--20 June 2019;pp.~7036--7045.
\item
  Liu, S.; Huang, D.; Wang, Y. Learning spatial fusion for single-shot
  object detection. arXiv 2019, arXiv:1911.09516.
\item
  Liu, S.; Qi, L.; Qin, H.; Shi, J.; Jia, J. Path aggregation network
  for instance segmentation. In Proceedings of the IEEE Conferenceon
  Computer Vision and Pattern Recognition (CVPR), Salt Lake City, UT,
  USA, 18--23 June 2018; pp.~8759--8768.
\item
  Tan, M.; Pang, R.; Le, Q.V. Efficientdet: Scalable and efficient
  object detection. In Proceedings of the IEEE/CVF Conference onComputer
  Vision and Pattern Recognition (CVPR), Seattle, WA, USA, 13--19 June
  2020; pp.~10781--10790.
\item
  Wang, C.Y.; Bochkovskiy, A.; Liao, H.Y.M. YOLOv7: Trainable
  bag-of-freebies sets new state-of-the-art for real-time
  objectdetectors. arXiv 2022, arXiv:2207.02696.
\item
  Xie, S.; Girshick, R.; Dollar, P.; Tu, Z.; He, K. Aggregated Residual
  Transformations for Deep Neural Networks. In Proceedingsof the 2017
  IEEE Conference on Computer Vision and Pattern Recognition (CVPR),
  Honolulu, HI, USA, 21--26 July 2017;pp.~1492--1500.
\item
  Zhou, T.; Zhao, Y.; Wu, J. Resnext and res2net structures for speaker
  verification. In Proceedings of the 2021 IEEE SpokenLanguage
  Technology Workshop (SLT), Shenzhen, China, 19--22 January 2021;
  pp.~301--307.
\item
  Neubeck, A.; Van Gool, L. Efficient non-maximum suppression. In
  Proceedings of the 18th International Conference on PatternRecognition
  (ICPR'06), Hong Kong, China, 20--24 August 2006; Volume 3,
  pp.~850--855.
\item
  Wang, C.Y.; Liao, H.Y.M.; Wu, Y.H.; Chen, P.Y.; Hsieh, J.W.; Yeh, I.H.
  CSPNet: A new backbone that can enhance learningcapability of CNN. In
  Proceedings of the IEEE/CVF Conference on Computer Vision and Pattern
  Recognition (CVPR), Seattle,WA, USA, 13--19 June 2020; pp.~390--391.
\item
  Yang, S.; Luo, P.; Loy, C.C.; Tang, X. Wider face: A face detection
  benchmark. In Proceedings of the IEEE Conference on ComputerVision and
  Pattern Recognition (CVPR), Las Vegas, NV, USA, 27--30 June 2016;
  pp.~5525--5533.
\item
  Rezatofighi, H.; Tsoi, N.; Gwak, J.; Sadeghian, A.; Reid, I.;
  Savarese, S. Generalized intersection over union: A metric and a
  lossfor bounding box regression. In Proceedings of the IEEE/CVF
  Conference on Computer Vision and Pattern Recognition (CVPR),Long
  Beach, CA, USA, 16--20 June 2019; pp.~658--666.
\item
  Yu, X.; Gong, Y.; Jiang, N.; Ye, Q.; Han, Z. Scale match for tiny
  person detection. In Proceedings of the IEEE/CVF WinterConference on
  Applications of Computer Vision (WACV), Snowmass Village, CO, USA,
  1--5 March 2020; pp.~1257--1265.
\item
  Zhu, P.; Wen, L.; Du, D.; Bian, X.; Fan, H.; Hu, Q.; Ling, H.
  Detection and Tracking Meet Drones Challenge. IEEE Trans. PatternAnal.
  Mach. Intell. 2021. {[}CrossRef{]} {[}PubMed{]}
\item
  Everingham, M.; Van Gool, L.; Williams, C.K.I.; Winn, J.; Zisserman,
  A. The PASCAL Visual Object Classes Challenge 2012(VOC2012) Results.
  Available online:
  http://www.pascal-network.org/challenges/VOC/voc2012/workshop/index.html(accessed
  on 30 July 2022).
\item
  Tong, K.; Wu, Y.; Zhou, F. Recent advances in small object detection
  based on deep learning: A review. Image Vis. Comput. 2020,97, 103910.
  {[}CrossRef{]}
\end{enumerate}

\end{document}
